\chapter{Intrauterine Device (IUD) Insertion}

\section*{Overview}
An IUD is a small, T-shaped contraceptive device inserted into the uterus. It prevents pregnancy by releasing copper (non-hormonal) or levonorgestrel (hormonal), altering the uterine environment.

\section*{Alternative Names}
IUD placement, Coil insertion

\section*{Principle}
A small T-shaped device is loaded into an insertion tube, passed through the cervix into the uterine cavity, and released so its arms open and it seats at the fundus. It prevents pregnancy through local inflammation, thickened cervical mucus, and, for hormonal devices, progestin release.
\section*{History}
Early intrauterine devices were developed in the early 20th century. Plastic, T-shaped copper IUDs were developed in the 1960s, and hormonal IUDs in the 1990s.

\section*{Why This Test or Procedure Is Ordered}
Ordered by patients seeking highly effective, long-acting reversible contraception (LARC) or management of heavy menstrual bleeding (hormonal IUD).

\section*{Is This a Primary or Secondary Test or Procedure}
Primary, highly effective method of long-acting reversible contraception.

\section*{How This Test or Procedure Is Performed}
The clinician performs a bimanual exam, inserts a speculum, stabilizes the cervix with a tenaculum, measures uterine depth, and inserts the IUD through the cervix using a sterile applicator tube.

\section*{What Preparation Is Needed}
A negative pregnancy test is mandatory. Taking an NSAID (ibuprofen) 1 hour before can help reduce insertion cramping.

\section*{Contraindications and Precautions}
Pregnancy, active pelvic infection (PID, mucopurulent cervicitis), unexplained vaginal bleeding, or uterine anomalies.

\section*{Risks}
Severe uterine cramping during insertion, vasovagal reaction, expulsion (<5%), pelvic infection (low risk, mostly in first 20 days), or uterine perforation (rare, <1 in 1,000).

\section*{What Happens After the Test or Procedure}
Cramping and irregular spotting are common for several weeks. Check the IUD strings monthly. Avoid tampon use for the first 24-48 hours.

\section*{Understanding Your Results}
Proper placement is confirmed if the strings are visible at the cervical os and the device is confirmed intrauterine on ultrasound (if checked).

\section*{Normal Results}
IUD positioned in the fundus of the uterus; strings visible at the external cervical os; no signs of displacement.

\section*{Follow-up Tests and Procedures}
String check in 4-6 weeks.

\section*{References}
\begin{enumerate}
  \item MedlinePlus. Intrauterine Device (IUD) Insertion. U.S. National Library of Medicine; 2025.
  \item Current Medical Diagnosis and Treatment. McGraw-Hill; 2025.
\end{enumerate}

\clearpage
